In this section we prove some basic properties about the girth of simple closed curves on the boundary of a handlebody.  For instance, we show that there are simple closed curves of arbitrarily high girth on the boundary of any handlebody.  We also introduce a conjecture equating the girth of a simple 
closed curve to another simple quantity.

Let $H$ be a handlebody.  Let $\gamma$ be a simple closed curve in $\bd H$.  
The main result of this section is the following Lemma~\ref{one embedded curve}.  When $\gamma$ is a simple closed curve, the minimum geometric intersection number of a curve homotopic to $\gamma$ and an essential disc in $H$ can be achieved by keeping $\gamma$ fixed.  In other words, no homotopy of $\gamma$ is necessary to realize the girth.

\begin{lemma}
\label{one embedded curve}
\label{simple closed curve}
Let $H$ be a handlebody.  Let $\gamma$ be a simple closed curve in $\bd H$.  Let $D$ be a girth-realizing disc for $\gamma$.  There exists an essential disc $D^{\prime}$ in $H$ such that $|\gamma\cap D^{\prime}|$ equals $\gamma\wedge D$ .  Equivalently, the girth of $\gamma$ equals the minimum geometric intersection number of $\gamma$ and an essential disc in $H$.
\end{lemma}

We give two proofs of Lemma~\ref{one embedded curve} in this section.  The first is somewhat shorter, but the construction is less direct.  It uses the tower construction in the proof of Dehn's Lemma.  The second proof is longer.  The construction is more explicit and easier to carry out.  We defer that proof to the end of this section.  There is even a third proof of Lemma~\ref{one embedded curve}.  The third proof can be obtained by following the proof of our main result, Theorem~\ref{subsurface theorem} in Section~\ref{maintheorem}.  That proof is of course much longer but entirely different in spirit.

\begin{proof}[First proof of Lemma~\ref{one embedded curve}.]
Let $n$ be the girth of $\gamma$.  The disc $D$ in $H$ is a girth-realizing disc for $\gamma$.  We will construct an essential disc $D^{\prime}$ in $H$ that intersects $\gamma$ exactly $n$ times.  Isotope $D$ in $H$ so that $D$ is transverse to $\gamma$.  Choose a single lift of $D$ to the universal cover $\pi:\hh\rightarrow H$.  The disc $D$ in the universal cover $\hh$ defines an interface $\beta$ for $\hh$.  By the definition of girth, the number of lifts of $\gamma$ that cross $\beta$ equals $n$. 

Let $B$ denote the boundary $\bd D$ of this lift to $\hh$.  The simple closed curve $B$ represents a nontrivial element of $H_{1}(\bd \hh,\Z)$, where $H_{i}$ denotes the usual homology.  We will construct a disjoint collection of simple closed curves in $\bd \hh$ that is homologous to $B$ and intersects $\pi^{-1}(\gamma)$ exactly $n$ times.  

Note that every curve in $\pi^{-1}(\gamma)$ can be homotoped to intersect $B$ either zero or one times.  The curves that cross $\beta$ can be made to intersect $B$ once, while the curves that do not cross $\beta$ can be made to intersect $B$ not at all.

View $B$ as a collection of curves.  Consider the lifts of $\gamma$ that intersect $B$ geometrically.  Since $|D\cap\gamma|$ is finite, there are finitely many such lifts.  For each such lift, modify $B$ in a small neighborhood so that the lift intersects the new $B$ either zero or one time.  This can be done while keeping the homology class of $B$ the same.  Lifts that cross $\beta$ will intersect $B$ once, while lifts that do not cross $\beta$ will be disjoint from $\beta$.  See Figure~\ref{curvemodify} for an illustration of how this can be done.

\begin{figure}
\epsfig{file=figures/curvemodify.eps, height=1.5in}
\caption{
This lift of $\gamma$ to the universal cover $\hh$ intersects the curve collection $B$ five times.  Since the lift crosses the interface $\beta$ once, the curve $B$ can be modifed to intersect the lift just once.  The change takes place in a neighborhood of the lift.  The new $B$ is homologous to the original.
}
\label{curvemodify}
\end{figure}

We now have a curve collection $B$ that intersects $\pi^{-1}(\gamma)$ exactly $n$ times.  Choose a curve component $B^{\prime}$ of $B$ that represents a nontrivial element of $H_{1}(\bd \hh,\Z)$.  This is possible since the original $B$ represented a nontrivial homology class.  This curve bounds a disc $D^{\prime}$ in $\hh$.

Project the disc $D^{\prime}$ to $H$.  The disc $D^{\prime}$ intersects $\gamma$ no more than $n$ times in $H$.  However, $D^{\prime}$ is not necessarily embedded.  To complete the proof, we apply the following Lemma~\ref{tower} to $D^{\prime}$.
\end{proof}

Lemma~\ref{tower} is where we use the tower construction.  Our other proof of Lemma~\ref{one embedded curve} is different.  In that proof, more care is taken in modifying $B$.  The modification is done in an equivariant way.  In all stages of that construction, the curves $B$ project down to a disjoint collection of simple closed curves in $H$.

\begin{lemma}
\label{tower}
Let $H$ be a handlebody.  Let $\gamma$ be a simple closed curve in $\bd H$.  Let $D$ be a disc in $H$ whose boundary is essential in $\bd H$ but is not necessarily embedded.  Then there is an essential embedded disc $D^{\prime}$ in $H$ with $|D^{\prime}\cap\gamma|\leq|D\cap\gamma|$.
\end{lemma}

For this proof, we assume familiarity with the tower construction.  In the course of the proof, observe that Lemma~\ref{tower} continues to hold true if the handlebody $H$ is replaced by any compact 3-manifold.  Moreover, if $\bd H$ has a metric, the result holds true if $|D\cap\gamma|$ is replaced by measuring the length of $\bd D$.  Note also that $\gamma$ need not be simple nor connected for the proof to hold.  For all of these variations, the same argument works.

\begin{proof}[Proof of Lemma~\ref{tower}.]
Assume that $\bd D$ is not a simple closed curve.  Homotope $\gamma$ in $\bd H$ so that $\gamma$ is transverse to $D$ and so that $\gamma\cap D$ lies away from the self-intersections of $\bd D$.  This can be done without increasing the intersection number $|\gamma\cap D|$.  Also make the self-intersections of $\bd D$ transverse, again without increasing $|\gamma\cap D|$.

Since we will be modifying $D$, view $D$ as a map $D^{2}\rightarrow D_{0}\subset H$.  Consider a regular neighborhood $V_{0}$ of $D_{0}$ in $H$.  In the tower construction, one considers a finite sequence of 2-sheeted covers, starting with a 2-sheeted cover of $V_{0}$.  At each stage of the construction, we lift the map of $D^{2}$.  At the top of the tower construction, we have a space $V_{n}$ and a map $D^{2}\rightarrow D_{n}\subset V_{n}$.  The space $V_{n}$ has no 2-sheeted cover, so $\bd V_{n}$ consists of 2-spheres.

Let $S^{2}$ denote the component of $\bd V_{n}$ that contains $\bd D_{n}$.  Consider the preimage of $\bd H$ in $S^{2}$ under the composition of the covering maps.  This is a compact subsurface of $S^{2}$ that contains $\bd D_{n}$.  Call this compact subsurface $S$.

At this point we need a simple closed curve $C$ in $S$ that projects to an essential closed curve in $\bd H$ and does not increase intersection with $\gamma$.  There are a couple ways to do this.  We choose to apply Lemma~\ref{top of tower}.  Rather than interrupt the flow of the proof, we prove this lemma afterwards.

% Choose a regular neighborhood of $\bd D_{n}$ inside $S$.  Let $S^{\prime}$ denote the closure of this regular neighborhood.  Away from the self-intersections of $\bd D_{n}$, the surface $S^{\prime}$ looks like an I-bundle.  The $1/2$ cross-section corresponds to $\bd D_{n}$, and the $0$ and $1$ cross-sections correspond to boundary circles of $S^{\prime}$.  Since $\bd D_{n}$ lies inside a 2-sphere, the $0$ and $1$ cross-sections at a given slice must correspond to different components of $\bd S^{\prime}$.

% Since $S^{\prime}$ contains the curve $\bd D_{n}$, the fundamental group $\pi_{1}(S^{\prime})$ projects to a nontrivial subgroup of $\pi_{1}(\bd H)$.  Also, since $S^{\prime}$ is contained in the 2-sphere $S^{2}$, the surface $S^{\prime}$ is a planar surface.  Thus $\pi_{1}(S^{\prime})$ is generated by the boundary curves of $S^{\prime}$.  In particular, at least one of its boundary curves must project to a nontrivial curve in $\bd H$.  Choose such a boundary curve $C$.

% By taking the regular neighborhood $S^{\prime}$ of $\bd D_{n}$ small enough, we ensure that $C$ parallels $\bd D_{n}$ closely.  Thus, where $\pi$ is the composition of the projection maps down to $H$, we have $|\pi(C)\cap \gamma|<|D\cap\gamma|$.  This is by the I-bundle remark above and because $D\cap\gamma$ was taken to lie away from the self-intersections of $\bd D$.

Apply Lemma~\ref{top of tower} to $\bd D_{n}$ mapping to $\pi_{1}(\bd H)$.  We obtain a simple closed curve $C$ in $S$ that projects to an essential closed curve in $\bd H$.  Since $C$ is a subset of $\gamma$, it does not increase intersection with $\gamma$ downstairs.  Thus, where $\pi$ is the composition of the projection maps down to $H$, we have $|\pi(C)\cap \gamma|\leq|D\cap\gamma|$.

Since the simple closed curve $C$ lies inside the 2-sphere $S^{2}$, it bounds an embedded disc $B_{n}$ inside $V_{n}$.  For the remainder of the proof, we descend through our tower of 2-sheeted covers, modifying $B_{n}$ at each stage.

At each stage, map the embedded disc $B_{i}$ down to a disc $B_{i-1}$ via the corresponding covering map.  Since the cover is 2-sheeted, $B_{i-1}$ has at most double points.  These can be resolved in the standard way by cutting-and-pasting.  This cutting-and-pasting does not increase intersection with $\gamma$ in $H$ because $\bd B_{i-1}$ is affected only near the preimage of self-intersections of $\bd D$.  The intersection points $D\cap\gamma$ were taken to be away from the self-intersections of $\bd D$.  In the end, we have an embedded disc $B_{0}$ in $H$.  Take $D^{\prime}$ to be $B_{0}$.  $D^{\prime}$ is an embedded essential disc in $H$ satisfying $|D^{\prime}\cap\gamma|\leq|D\cap\gamma|$.
\end{proof}

We now prove the final lemma needed for the proof of Lemma~\ref{one embedded curve}.  Lemma~\ref{top of tower} is needed at the top of the tower construction in the proof of Lemma~\ref{tower}.

\begin{lemma}
\label{top of tower}
Let $S$ be a surface.  Let $G$ be a group.  Let $\Phi:\pi_{1}(S)\rightarrow G$ be a homomorphism from the fundamental group of $S$ to $G$.  Let $\gamma$ be a curve in $S$ with transverse self-intersections.  Suppose $\Phi$ maps $\gamma$ to a nontrivial conjugacy class in $G$.  Then $\gamma\subset S$ contains a simple closed curve that maps to a nontrivial conjugacy class in $G$.
\end{lemma}

\begin{proof}
The proof is by downward induction on the number of self-intersections of $\gamma$.  Choose a point in $\gamma$ and begin tracing out the curve $\gamma$ in one direction.  Stop when the curve first intersects itself.

After dropping the initial arc segment, this forms a simple closed curve $\gamma^{\prime}$ in $S$.  if $\gamma^{\prime}$ maps to a nontrivial conjugacy class in $G$, we are done.  If $\gamma^{\prime}$ maps trivially to $G$, then delete $\gamma^{\prime}$ from $\gamma$ to form a new curve $\gamma$.  See Figure~\ref{selfintersection} for an illustration.  The homomorphism $\Phi$ still maps the new $\gamma$ to a nontrivial conjugacy class.

Repeat the construction until the proof is complete.
\end{proof}

\begin{figure}
\epsfig{file=figures/selfintersection.eps, height=2in}
\caption{
We show part of a curve $\gamma$ in a surface $S$.  The curve $\gamma$ maps to a nontrivial conjugacy class in a group $G$.  If the loop on the left maps trivially, then removing it creates a curve with fewer self-intersections that continues to map nontrivially.
}
\label{selfintersection}
\end{figure}

We now return to our main narrative.  Let $\pi:\hh\rightarrow H$ denote the universal covering space of $H$.  Define a graph $\Gamma$ in $\bd \hh$ as follows.  Take the components of $\bd\hh-\pi^{-1}(\gamma)$ to be the vertices of $\Gamma$.  Take the components of $\pi^{-1}(\gamma)$ to be the edges.  Each component of $\pi^{-1}(\gamma)$ is a curve that borders two components of $\bd\hh-\pi^{-1}(\gamma)$, which are possibly the same.  In this way, each edge defines two vertices.  We call $\Gamma$ the \emph{graph associated} to $\gamma$.  The graph $\Gamma$ is locally infinite if the genus of $H$ is two or more.

The girth of a graph is the length of a shortest cycle in the graph.  By the following lemma, the girth of a simple closed curve $\gamma$ equals the girth of the graph $\Gamma$ associated to $\gamma$.  This explains our choice of the word girth.  Lemma~\ref{graph girth} is a corollary of Lemmas~\ref{one embedded curve} and~\ref{tower}.

\begin{lemma}
\label{graph girth}
Let $H$ be a handlebody.  Let $\gamma$ be a simple closed curve in $\bd H$.  Let $\Gamma$ be the graph associated to $\gamma$.  If $\bd H-\gamma$ is incompressible in $H$, then the girth of $\gamma$ equals the girth of $\Gamma$.
\end{lemma}

By Lemma~\ref{one embedded curve}, the surface $\bd H-\gamma$ is compressible if and only if $\gamma$ has zero girth.  

\begin{proof}
Let $n$ be the girth of $\gamma$.  By Lemma~\ref{one embedded curve}, the smallest intersection number of $\gamma$ with an essential disc in $H$ equals $n$.  Let $D$ be an essential embedded disc that achieves this minimum.  Then $D$ intersects $\gamma$ in $n$ points.  Since $\bd H-\gamma$ is incompressible, $n$ is bigger than zero.

Lift $D$ to the universal cover $\hh$.  Since $n$ is bigger than zero, the boundary of $D$ defines a circuit $C$ in the graph $\Gamma$ associated to $\gamma$.  The circuit $C$ has $n$ edges.  Thus, the girth of $\Gamma$ is less than or equal to the girth of $\gamma$.

We now prove the reverse inequality.  Let $m$ be the girth of the graph $\Gamma$ in the graph theory sense.  Let $C$ be a circuit in $\Gamma$ that achieves the girth of $\Gamma$.  Then the length of $C$ is $m$.  Realize $C$ as the boundary of a disc $D$ in $\hh$.  Then $D$ intersects lifts of $\gamma$ exactly $m$ times. Map $D$ to $H$.  Then $D$ is a disc in $H$ that intersects $\gamma$ exactly $m$ times but is not necessarily embedded.

Apply Lemma~\ref{tower} to the non-embedded disc $D$ to create an embedded disc $D$ that intersects $\gamma$ no more than $m$ times.  
The girth of $\gamma$ equals the minimum geometric intersection number of an essential embedded disc in $H$ and a curve homotopic to $\gamma$.  Thus the girth of $\gamma$ is less than or equal to $m$, the girth of $\Gamma$.
\end{proof}

In the case that $\gamma$ is simple, we conjecture another interpretation of girth.  We conjecture that the girth of $\gamma$ equals the minimum geometric intersection number of $\gamma$ and a curve in $\bd H$ that is homotopic to $\gamma$ in $H$ but not in $\bd H$.  In other words, the girth of $\gamma$ could be computed by ranging over all curves homotopic to $\gamma$ rather than over all essential discs in $H$.

\begin{conjecture}
\label{special conjecture}
Let $H$ be a handlebody.  Let $\gamma$ be a simple closed curve 
in $\bd H$.  Let $\kappa$ be a closed curve in $\bd H$ 
that is homotopic to $\gamma$ in $H$ but not isotopic to 
$\gamma$ in $\bd H$.  Then the intersection number of $\kappa$ and 
$\gamma$ is at least the girth of $\gamma$.
\end{conjecture}

To obtain a curve $\kappa$ homotopic to $\gamma$ such that $|\kappa \cap \gamma|$ equals the girth of $\gamma$, simply Dehn-twist $\gamma$ around a girth-realizing disc for $\gamma$. 

By Lemma~\ref{one embedded curve}, a simple closed curve of girth zero is a curve disjoint from some essential disc in $H$.  By the following lemma, simple closed curves of girth one do not exist in the higher genus case.  The argument is the classical argument that stabilized Heegaard splittings of genus bigger than one are reducible.  We repeat it here for completeness.
  
\begin{lemma}
\label{geometric girth one}
Let $H$ be a handlebody of genus two or more.  Let $\gamma$ be a closed curve in $\partial H$.  If $\gamma$ is simple, then $\gamma$ does not have girth one.
\end{lemma}

\begin{proof}
Let $\gamma$ be a curve in $\bd H$ satisfying $\girth(\gamma)\leq 1$ in $H$.
By Lemma~\ref{simple closed curve}, there is an essential disc $D$ that
intersects $\gamma$ either once or not at all.  Suppose a disc $D$ in $H$ transversely intersects $\gamma$ exactly once. Consider the boundary of a regular neighborhood in $\bd H$ of $\gamma\cup\bd D$.
This curve bounds a disc in $H$ and intersects $\gamma$ zero times.
Since $H$ has genus bigger than one, the curve is essential in $\bd H$.
Thus the girth of $\gamma$ is zero.
\end{proof}

Lemma~\ref{geometric girth one} is actually true for all curves $\gamma$ in $H$, but the proof is more difficult.  We prove this as Lemma~\ref{group girth one}.  For now we prove the following special case.  Lemma~\ref{genus two} is a corollary of Lemma~\ref{geometric girth one}.

\begin{lemma}
\label{genus two}
Let $H$ be a handlebody.  Let $\gamma$ be a closed curve in $H$.  If the genus of $H$ equals two, then $\gamma$ does not have girth one.
\end{lemma}

\begin{proof}
Let $\gamma$ be a closed curve in $H$ that intersects an essential disc $D$ exactly once.  It suffices to show that $\gamma$ has girth zero.  Extend $D$ to a disc system $\dd$ for $H$.  Let $\{a, b\}$ be the corresponding generating system for $\pi_{1}(H)$, with $a$ corresponding to $D$.  The curve $\gamma$ is conjugate to a group element of the form $a b^{m}$, for some integer $m$.  Homotope $\gamma$ to a simple closed curve in $\bd H$ that intersects $D$ exactly once.  By Lemma~\ref{geometric girth one}, the girth of $\gamma$ must be zero.
\end{proof}

Whereas Lemma~\ref{geometric girth one} shows that simple closed curves of girth one do not exist, Lemma~\ref{large girth} shows that simple closed curves of large girth do exist.  In particular, every free group contains elements of arbitrarily high girth.  

\begin{lemma}
\label{large girth}
Let $H$ be a handlebody.  The surface $\bd H$ contains simple closed
curves of arbitrarily high girth in $H$.
\end{lemma}

\begin{proof}
Let $g$ be the genus of $H$.  Let $S$ be a closed disc with $g$ open discs removed.  Then $H$ is homeomorphic to the product of $S$ with a closed
interval.  Identify $H$ with $S\times I$.  The boundary of the cross-section $S\times\{1/2\}$ is a collection of $g+1$ closed curves in $\bd H$.  Denote this collection $\gamma$. 

The collection of curves $\gamma$ represents a rational element of $\ps(\bd H)$, the projective space of measured laminations on $\bd H$.  See~\cite{Thu2} an understanding of $\ps(\bd H)$.  Perturb $\gamma$ to an irrational element $\lambda$ in $\ps(\bd H)$.  The simple closed curves in $\bd H$ are dense in $\ps(\bd H)$.  By approximating $\lambda$ by simple closed curves, we can generate curves $\gamma^\prime$ of higher and higher girth.

Consider a standard collection of $3g-3$ non-parallel essential discs in $H$, so that each disc intersects $\gamma$ twice.  See Figure~\ref{highgirth}.  Call this collection $\dd$.  The collection $\dd$ divides $H$ into $2g-2$ filled-in pairs of pants.  Each pair of pants has three boundary arcs that connect the three disc boundaries in a cycle.

\begin{figure}
\epsfig{file=figures/highgirth.eps, height=1.25in}
\caption{
The thin curves are the boundaries of $3g-3$ discs in a handlebody of genus $g$.  Each disc intersects the collection $\gamma$ of thick curves in exactly two points.
}
\label{highgirth}
\end{figure}

Consider the measured lamination $\lambda$ with respect to $\dd$.
If the perturbation is small enough, the lamination 
$\lambda$ will intersect each pair of pants in three bands of arcs 
in the same isotopy class as the original arcs.  See 
Figure~\ref{highgirthpant}.  Each band contains an infinite number 
of arcs, and each band has roughly the same measure.  

\begin{figure}
\epsfig{file=figures/highgirthpant.eps, height=1.25in}
\caption{
Perturb the curve $\gamma$ of Figure~\ref{highgirth}  
to an irrational measured lamination $\lambda$.
The lamination $\lambda$ intersects each of the $2g-2$ pants 
in three bands of arcs, after isotopy.}
\label{highgirthpant}
\end{figure}

Let $\gamma^\prime$ be a simple curve that approximates $\lambda$.
The curve $\gamma^\prime$ must intersect each pair of pants in 
three bands of arcs in the same isotopy class as the bands for 
$\lambda$.  Each band for $\gamma^\prime$ has roughly the same 
number of arcs.  Let $n$ denote the smallest number of these 
arcs among all $6g-6$ bands of arcs.  As $\gamma^\prime$ more 
closely approximates $\lambda$, the number $n$ must grow larger.

Let $D$ be an essential disc in $H$ that minimizes intersection with $\gamma$.  By Lemma~\ref{one embedded curve}, the geometric intersection number of $D$ and $\gamma$ equals the girth of $\gamma$.  It is evident that $D$ must intersect $\gamma^\prime$ in at least $2n$ points.  See Figure~\ref{highgirthpant} for an illustration.  Thus the girth of $\gamma$ is at least $2n$.
\end{proof}

We now present our second proof of Lemma~\ref{one embedded curve}.

\begin{proof}[Second proof of Lemma~\ref{one embedded curve}.]
Let $n$ be the girth of $\gamma$.  Recall that $\gamma$ is a simple closed curve in $\bd H$.  Let $D$ be a disc in $H$ that realizes the girth of $\gamma$.  We will construct a disc $D^{\prime}$ in $H$ that intersects $\gamma$ exactly $n$ times.

Let $\pi:\hh\to H$ be the universal cover.  Let $\Gamma$ be the tree associated to $D$.  The action of $\pi_{1}(H)$ on $\hh$ induces an action of $\pi_{1}(H)$ on $\Gamma$.  Note that the action of $\pi_{1}(H)$ on $\Gamma$ is not free when $\genus(H)\geq 2$.  Isotope $\gamma$ in $\bd H$ to be transverse to $D$.  Lift $\gamma$ to a path $P$ in $\Gamma$.  The path $P$ is invariant under the cylic action of $\langle\gamma\rangle$ on $\Gamma$.

Descend to $\Gamma/\langle\gamma\rangle$.  Let $C$ denote the path $P/\langle\gamma\rangle$ in $\Gamma/\langle\gamma\rangle$.  The path $C$ is a map of a circuit to $\Gamma/\langle\gamma\rangle$.  The edges of $C$ are in correspondence with the points $\gamma\cap D$ in $H$, and the vertices of $C$ are in correspondence with the arcs $\gamma-D$ in $H$.

Consider a single lift of $D$ to $\hh$.  Let $B$ denote the boundary of this lift.  The curve $B$ is a simple closed curve that separates $\bd\hh$ into two noncompact pieces.  It represents a nontrivial element of $H_{1}(\bd \hh,\Z)$, where $H_{i}$ denotes the usual homology.  The curve $\pi(B)$ is the same as $\bd D$.  The points $\gamma\cap\pi(B)$ are in correspondence with $\gamma\cap D$.

Rename $\gamma$ as $\gamma^{\prime}$.  We reserve $\gamma$ for the original curve.  The proof proceeds by downward induction on the number of edges of $C$.    At each stage we modify the path $C$ in $\Gamma/\langle\gamma\rangle$, the curve $\gamma^{\prime}$ in $H$, and the curve $B$ in $\bd \hh$.  We stop when $C$ is a geodesic path in $\Gamma/\langle\gamma\rangle$.  Throughout the induction we preserve the correspondence between the edges of $C$, the points of $\gamma^{\prime}\cap D$, and the points of $B\cap \gamma$.  Note that after the first stage, the curve collection $B$ can have more than one connected component.  

\textbf{*Begin induction.}
Consider the path $C$.  If $C$ has no backtracking  in $\Gamma/\langle\gamma\rangle$, then skip to the end of the induction stage (\textbf{**}).  Otherwise, consider a backtracking of $C$.  The backtracking consists of a vertex and two edges of $C$.  The vertex of the backtracking corresponds to an arc $\kappa^{\prime}$ of $\gamma^{\prime}-D$ in $H$.

Homotope $\gamma^{\prime}$ in $H$ to eliminate the two points of intersection between the closure $\overline{\kappa^{\prime}}$ of $\kappa^{\prime}$ and the disc $D$.  Choose the homotopy so that $\gamma^{\prime}$ stays fixed outside a small regular neighborhood of $\overline{\kappa^{\prime}}$.  This reduces $|\gamma^{\prime}\cap D|$ by two and keeps all other points of $\gamma^{\prime}\cap D$ the same.  See Figure~\ref{homotopy} for an illustration.

\begin{figure}
\epsfig{file=figures/homotopy1.eps, height=1.75in}
\caption{
We show a homotopy of $\gamma^{\prime}$ near the arc $\kappa$.  In the beginning, the arc $\kappa$ connects the two intersection points of $\gamma^{\prime}$ and the shaded disc $D$.  The homotopy reduces $|\gamma^{\prime}\cap D|$ by two and keeps $\gamma^{\prime}$ fixed outside a small neighborhood of $\overline{\kappa}$.
}
\label{homotopy}
\end{figure}

Consider the path $C$ in $\Gamma/\langle\gamma\rangle$.  Collapse the backtracking in $C$ to create a new path in $\Gamma/\langle\gamma\rangle$ with two fewer edges.  The new path $C$ represents a lift of the new path $\gamma^{\prime}$ to $\Gamma/\langle\gamma\rangle$ .

We now modify the disjoint collection of simple closed curves $B$ in $\hh$.  See Figure~\ref{recursionstep} for an illustration of this construction.  The covering map $\pi$ embeds $B$ into $\bd H$.  Consider the two points of $\gamma^{\prime}\cap D$ corresponding to the backtracking we chose above.  These two points correspond to two points of $\gamma\cap D$, and they lie in $\pi(B)$.  Consider the arc $\kappa$ of $\gamma-D$ that connects these two points in $\bd H$ and is homotopically trivial in $H \rel D$.  By construction, the arc $\kappa$ satisfies the property that its closure $\overline{\kappa}$ intersects $\pi(B)$ only at the endpoints of $\overline{\kappa}$.

\begin{figure}
\epsfig{file=figures/recursionstep.eps, height=2in}
\caption{
The curve $\gamma$ intersects the thick curve $B$ on the left in a homotopically trivial arc $\kappa$.  We alter $B$ by considering a thin neighborhood $U$ around the arc $\kappa$, shown in light gray.  Consider a pair of arcs in $U$ running alongside $\kappa$.  Now cut and paste to form the new collection $B$.  The new collection is homologous to the old because the four arcs bound the rectangular region shown in darker gray.
}
\label{recursionstep}
\end{figure}

Let $U$ in $\bd H$ be a neighborhood of $\kappa$ chosen small enough so that (1) $U$ intersects $\pi(B)$ only in a small arc around each endpoint of $\overline{\kappa}$, and (2) the intersection $U\cap\gamma$ is a single arc component and a small neighborhood of $\kappa$ in $\gamma$.  Lift this picture back to $\hh$ with $\pi^{-1}$.  Replace the two arcs of $B\cap U$ in $U$ with a homologous pair of arcs that run parallel to $\kappa$. Choose the pair arcs so they do not intersect $\pi^{-1}(\gamma)$.

This gives us a new disjoint collection of simple closed curves $B$ in $\bd \hh$.  The collection $B$ satisfies the following three properties.  The new collection $B$ represents the same element of $H_{1}(\bd \hh, \Z)$ as before.  The map $\pi$ embeds $B$ into $H$.  The collection $\pi(B)$ intersects $\gamma$ in two fewer points than before.

Repeat the induction stage by returning to (\textbf{*}).

\textbf{**End induction stage.}

We now have a path $C$ in $\Gamma/\langle\gamma\rangle$ without backtrackings, a disjoint collection of simple closed curves $B$ in $\bd\hh$, and a curve $\gamma^{\prime}$ in $H$ homotopic to $\gamma$.  The edges of $C$ in $\Gamma/\langle\gamma\rangle$, the points of $\gamma\cap\pi(B)$ in $\hh$, and the points of $\gamma^{\prime}\cap D$ in $H$ are all in correspondence.

Lift $\gamma^{\prime}$ to a path $P^{\prime}$ in $\Gamma$.  The path $C$ is the projection of $P^{\prime}$ to $\Gamma/\langle\gamma\rangle$.  Since $C$ has no backtrackings, the path $P^{\prime}$ has no backtrackings.  By Lemma~\ref{geodesic}, the curve $\gamma^{\prime}$ minimizes intersection with $D$ over its homotopy class.  By Lemma~\ref{wedge}, the intersection number $|\gamma^{\prime}\cap D|$ equals $\girth(\gamma^{\prime})$.  Since $\gamma^{\prime}$ is homotopic to $\gamma$, the intersection number equals $\girth(\gamma)$.  Hence $C$ has $n$ edges.  Since the edges of $C$ are in correspondence with $\gamma\cap\pi(B)$, the curve $\gamma$ intersects $\pi(B)$ in $n$ points.

Since $B$ represents a nontrivial homology class in $\bd \hh$, at least one component of $B$ represents a nontrivial homology class.  Choose one of these components of $B$, and call it $B^{\prime}$.  Map $B^{\prime}$ to $H$ with $\pi$.  The curve $\pi(B^\prime)$ is a simple closed curve in $\bd H$ that bounds an embedded disc $D^\prime$ in $H$ and intersects $\gamma$ in no more than $n$ points.

The disc $D^{\prime}$ is essential because $B^{\prime}$ represents a nontrivial homology class in $\bd \hh$.  The disc $D^\prime$ intersects $\gamma$ in no more than $n$ points.  Since the girth of $\gamma$ is $n$, the disc $D^\prime$ intersects $\gamma$ in at least $n$ points.  The disc $D^{\prime}$ intersects $\gamma$ in exactly $n$ points, so we are done.
\end{proof}

The proof above constructs a geometrically girth-realizing disc $D^{\prime}$ from an algebraically girth-realizing disc $D$.  The construction differs depending on the order in which the backtrackings of $C$ were canceled.  In general, different sequences are possible.  Figure~\ref{parenthesize} shows a simple example with two possibilities.  We encourage the reader to construct the two resulting curve collections $B$.

\begin{figure}
\epsfig{file=figures/parenthesize.eps, height=2in}
\caption{
The curve $\gamma$ intersects the shaded disc $D$ in four points.  The corresponding lift of $\gamma$ to $\Gamma$ has three backtrackings.  We illustrate two orders in which the backtrackings can be canceled, as in our second proof of Lemma~\ref{one embedded curve}.  The arrows depict the first cancellations.}
\label{parenthesize}
\end{figure}
